\chapter{Discussion}\label{Discussion}

\section{Limitations}

The implemented architecture was evaluated as a hardware realization of the proposed \ac{SBPS}+\ac{SE} carrier-recovery concept. The most important configurable parameters in this evaluation are the number of \ac{SBPS} refinement stages, the number of radius-directed stages, the processing block size and the number of phase-compensation sections. These parameters correspond directly to the algorithmic assumptions introduced in Chapter~\ref{Proposed phase recovery architecture}: \ac{SBPS} is responsible for a sufficiently accurate centre-phase estimate, while \ac{SE} and section-based compensation assume that the remaining intra-block phase evolution can be represented by a signed spread value around this centre phase.

The first limitation is therefore the locally linear phase assumption from Equation~\ref{eq:SE-linear-assumption}. This assumption is introduced in the system overview of Chapter~\ref{Proposed phase recovery architecture} and then used explicitly in the spread-estimation derivation in Section~\ref{input selector} and in the compensation model of Equations~\ref{eq:section-phase}--\ref{eq:section-residual-error}. The results confirm both the usefulness and the boundary of this assumption. Figure~\ref{fig:results-evm-sections-cfo} shows that increasing the number of compensation sections improves the result when the residual phase evolution is mainly a deterministic slope. However, Figure~\ref{fig:results-residual-EVM-by-symbol_nondefault} shows that under a more demanding residual-\ac{CFO} case, the remaining error stays position-dependent even with more compensation granularity. This indicates that spread-estimation accuracy, unwrap robustness and non-linear linewidth-induced phase variation also become limiting factors.

The second limitation is related to phase unwrapping. The unwrapper assumes that the leading and lagging edge estimates of neighbouring blocks provide a reliable continuity test, as formalized in Equation~\ref{eq:unwrap-edge-diff}. Its quadrant update in Equation~\ref{eq:unwrap-index-update} is therefore reliable only as long as the estimated edge discontinuity remains within the \(\pm\pi/4\) decision margin. A poor spread estimate can affect the system twice: it first causes an inaccurate intra-block compensation profile, and it can then move the edge-continuity decision closer to a wrong branch. In mild cases this appears as increased edge \ac{EVM}; in severe cases it can become a cycle slip.

A third limitation follows from the fixed input scaling. Chapter~\ref{Hardware Implementation} explains the Q1.15 sample representation and the static constellation scaling used to preserve headroom during rotations and moderate noise. Figures~\ref{fig:clipping_false_minima} and~\ref{fig:clipping_constellation} show why this is necessary: clipping can distort the decision-directed \ac{SBPS} metric and create false minima. The same scaling choice also affects the fourth-power path, because saturated samples would bias the selected-symbol and dot-product calculations. The implemented receiver is therefore not an adaptive front-end; changing the input dynamic range or operating at substantially different \ac{SNR} conditions would require renewed fixed-point validation.

The fourth limitation is the dependence of \ac{SE} on suitable \ac{16QAM} symbols. The input-symbol considerations in Section~\ref{input selector} explain why middle-ring symbols are less useful after the fourth-power operation and why the estimator selects representative candidates near the beginning, middle and end of a block. The hardware version of this selector is described in the \ac{SE} input-selector section of Chapter~\ref{Hardware Implementation}. This keeps the design compact, but it makes the spread estimate sequence-dependent. Under unfavourable symbol patterns the estimator may have to hold previous candidates or use a less reliable pair, which can reduce the quality of the signed spread estimate.

Finally, the \ac{RTL} contains several generics, but the complete verification flow is centred around the selected Q1.15 sample representation, Q4.28 phase representation and implemented lookup-table widths. These representations are summarized in Figure~\ref{fig:dataflow-top} and in Table~\ref{tab:top-design-streams}. Changing these widths affects clipping behaviour, accumulator growth, \ac{LUT} addressing, \ac{CORDIC} conversion and valid-signal alignment. The current results should therefore be interpreted as validation of the selected implementation point, not as proof that every parameterized variant is already verified.

\section{Implementation Trade-Offs}

The implementation deliberately favours explicit timing alignment over a more flexible streaming interface. The top-design overview in Chapter~\ref{Hardware Implementation} shows that the raw I/Q block is sent in parallel to \ac{SBPS}, \ac{SE} and an I/Q delay branch. The \textit{synchroniser} delays the signed spread estimate to match the \ac{SBPS} centre-phase path, while \textit{iq\_block\_delay} delays the raw samples until the corresponding phase estimates are available. This makes the timing contract easy to verify: when \textit{PUPC} receives a valid phase estimate, it also receives the spread estimate and sample block belonging to the same input block.

The cost of this explicit alignment is that the wrapper is not yet a fully elastic receiver interface. Most internal datapaths are pipelined, and the spread-estimation backend can accept aligned valid inputs once filled, as shown in Tables~\ref{tab:spread-estimation-latency} and~\ref{tab:sign-estimation-latency}. However, the complete design still relies on fixed valid pulses and deterministic delay lines rather than on a ready/valid protocol with backpressure. This is appropriate for the verification flow of Chapter~\ref{Verification and simulation}, but a production receiver would require a more general streaming integration.

The mixed phase representation is another important trade-off. The \ac{SBPS} stages use sine and cosine coefficients because the dominant operation is Cartesian rotation followed by slicing and metric accumulation, as described by Equation~\ref{eq:rotated_samples}. This is efficient for the \ac{SBPS} metric, but inconvenient for phase addition and interpolation. The downstream \ac{PU} and \ac{PC} blocks therefore use Q4.28 phase words, which make edge reconstruction, branch unwrapping and section-phase generation simple. The conversion cost at the boundary is justified because each representation is used where it best matches the required operation.

The resource results also show a practical timing trade-off. Table~\ref{tab:results-resource-usage} and Table~\ref{tab:results-timing-closure} show that the evaluated design is not limited by total area on the selected XCU200 device. Instead, the critical issue was a specific arithmetic path in the spread-estimation reciprocal-multiplication logic. Adding one local pipeline register allowed the design to meet the \(8~\mathrm{ns}\) timing target with negligible area increase. This illustrates that the main hardware challenge is not only reducing operation count, but also placing registers where wide fixed-point arithmetic would otherwise dominate the clock period.

\section{Possible Improvements}

One direct improvement is a more systematic fixed-point exploration. The present design validates one main numerical configuration. Sweeping the sample width, phase-word width, fourth-power output width, arccosine-\ac{LUT} address width and section-coefficient precision would show which parts of the design are over-provisioned and which parts are accuracy-limiting. The spread-estimation result in Figure~\ref{fig:se-relative error} makes this especially relevant near zero spread, where small normalized-dot-product errors are amplified by the nonlinear \(\arccos(\cdot)\) mapping.

The input selector could also be improved. The current design searches a limited number of symbols near the beginning, middle and end of the block, then relies on the candidate-valid logic and hold behaviour described in Chapter~\ref{Hardware Implementation}. A more robust version could search a larger window, rank candidates by distance from ambiguous decision regions, or use reliability information from the slicer. Such changes would mainly target cases where wrong spread signs or stale spread values create edge-localized residual error.

Another possible improvement is adaptive parameter selection. The results in Figures~\ref{fig:results-ber-stages-cfo} and~\ref{fig:results-evm-sections-cfo} show that the useful number of \ac{SBPS} stages and compensation sections depends on the channel condition. A receiver that estimates whether it is operating in a benign or demanding phase-noise regime could reduce switching activity by disabling unnecessary refinement stages or compensation sections when the full correction depth is not needed.

The \ac{SBPS} rotation branches are also candidates for alternative implementation. The current design uses multiplier-based rotations to keep the branch latency short and predictable. A pipelined \ac{CORDIC}-based rotator could reduce DSP usage, but it would increase latency and register usage and would require new top-level alignment settings. This option is therefore not automatically better, but it becomes relevant if DSP resources are more constrained than in the evaluated FPGA.

Finally, the architecture could be extended beyond the current \ac{16QAM} setting. Higher-order \ac{QAM} formats would require different candidate-selection rules, different radius normalization and possibly different modulation-removal logic. This extension would determine whether the same low-complexity centre-phase-plus-spread principle remains useful for coherent links with higher spectral efficiency.

\section{Societal Responsibility}

The societal relevance of this work comes from its intended application context: optical interconnects inside datacenters. As discussed in the introduction, datacenter traffic continues to grow because of cloud services, streaming platforms and AI workloads. Improving the efficiency of coherent receiver \ac{DSP} can therefore contribute to communication systems that provide higher throughput without increasing hardware complexity proportionally.

The direct stakeholders are the designers and operators of optical transceivers and datacenter networks, who must balance bandwidth, latency, power consumption and cost. Users of datacenter services are indirect stakeholders, because the scalability and reliability of these services depend on the interconnect fabric. The environment is also a stakeholder. Datacenters consume substantial electrical power, and optical transceivers contribute to this demand. A lower-complexity carrier-recovery backend can reduce part of the digital processing cost, but it should not be presented as a complete sustainability solution. Efficiency improvements reduce the energy per processed bit, while total environmental impact still depends on deployment scale, energy source and growth in demand.

For that reason, the value of this thesis lies in making one receiver function more hardware-efficient and more transparent. The design choices, fixed-point assumptions, verification flow and limitations are explicitly documented, allowing future work to compare accuracy, resources and timing instead of treating the receiver backend as a black box.

The same transparency also prevents overclaiming. The architecture is evaluated for a specific modulation format, block size and impairment range. It supports more efficient coherent interconnect research, but practical deployment would still require integration into a complete receiver, power measurements on the target technology and evaluation under broader operating conditions. Within this scope, the thesis contributes to responsible system design by reducing unnecessary processing complexity while clearly identifying the remaining technical and environmental caveats.
